\documentclass[10pt,a4paper]{amsart}
\usepackage[margin=19mm]{geometry}
\usepackage{amsmath,amssymb,mathtools}
\usepackage[scr=rsfs,cal=euler]{mathalfa}
\usepackage{xfrac}
\usepackage[unicode,hidelinks,bookmarks=false]{hyperref}
\newcommand{\ie}{i.e.\ }
\newcommand{\cf}{cf.\ }
\newcommand{\resp}{respectively\ }
\newcommand{\Subsection}{\subsection}
\providecommand{\nobreakdash}{\nobreak\hbox{-}}
\setlength{\emergencystretch}{2em}
\multlinegap0pt
\usepackage{amssymb}
\newtheorem{theorem}{Theorem}
\title{Exact values and bounds for Ramsey numbers of $C_4$ versus a star graph}
\author{Luis Boza}
\date{Source: arXiv:2409.12770v2 (CC BY 4.0)}
\begin{document}
\maketitle

\noindent\emph{Source and licence.} Exact values and bounds for Ramsey numbers of $C_4$ versus a star graph. Version arXiv:2409.12770v2 (CC BY 4.0), distributed by arXiv under the Creative Commons Attribution 4.0 International licence, http://creativecommons.org/licenses/by/4.0/, as recorded in the arXiv licence metadata for this exact version and shown on the abstract page https://arxiv.org/abs/2409.12770v2. Full licence notice: \texttt{licenses/arxiv-2409.12770-CC-BY-4.0.txt}.\par\smallskip

\noindent\textit{Editorial scope.} Counted unit: the article's theorem teo2 (source0.tex lines 121--123). The article's complete original proof of it is delivered with it (source0.tex lines 125--160, one \texttt{proof} environment of 36 lines). No mathematics is written, repaired or re-derived: the statement and the proof are the article's own lines, in the article's own notation, and no retained line is substituted or re-flowed.  No further source-local context is needed: every cross-reference of every retained line resolves inside the two delivered spans.  Nothing was omitted from any retained span, so the package carries no omission record.  No retained line contains a citation, so the wrapper carries no bibliography and no bibliography entry.  No retained line contains a figure environment, an image-inclusion command or drawing code, so no image is delivered and the package declares no asset dependency.  Editorial formatting only, in addition: an AMS \texttt{amsart} preamble that restates the article's own declarations and macros used by the retained lines, namely the environments \texttt{theorem} on separate counters, together with the standard packages those lines need.  The retained environments print in this document as Theorem 1 in order of appearance, whereas the article prints the same items with the numbering of its own full document.  The complete native source of the article as distributed in the arXiv e-print for this exact version is bundled unchanged as \texttt{sources/165508/source0.tex}.
\begin{theorem} \label{teo2}
Let $m\equiv 2 \,(\text{mod}\,\, 6)$ with $m\ge 8$. Then $f(m^2+3)\le m^2+m+4$.
\end{theorem}

\begin{proof}
Suppose there exists a graph $G$ with $m^2+m+4$ vertices such that
$G\nsupseteq C_4$ and $G\nsupseteq K_{1,m^2+3}$. Then $\delta(G)\ge m+1$.

Fix a vertex $v\in V(G)$. Since $G\nsupseteq C_4$, every vertex of $N_G(v)$ is adjacent
to at most one other vertex of $N_G(v)$; otherwise two such neighbours together with $v$
would form a copy of $C_4$. Moreover, if a vertex
$w\in V(G)\setminus (N_G(v)\cup\{v\})$ were adjacent to two vertices of $N_G(v)$, then
$v$, $w$, and these two vertices would also form a $C_4$. Hence every vertex of
$N_G(v)$ is adjacent to at least $d_G(v)-2$ vertices outside $N_G(v)\cup\{v\}$.

If $d_G(v)\ge m+2$, then there are at least
$(m+2)(m-1)$ edges joining $N_G(v)$ to
$V(G)\setminus (N_G(v)\cup\{v\})$. Since
$|V(G)\setminus (N_G(v)\cup\{v\})|
\le m^2+m+4-1-(m+2)<(m+2)(m-1)$,
some vertex outside $N_G(v)\cup\{v\}$ must be adjacent to at least two vertices of
$N_G(v)$, a contradiction. Therefore $d_G(v)\le m+1$.

Since $\delta(G)\ge m+1$, it follows that every vertex of $G$ has degree exactly $m+1$.

For each $v\in V(G)$, let $t_v$ denote the number of edges with both endpoints in $N_G(v)$. $t_v$ is the number of triangles to which $v$ belongs, and the total number of triangles in $G$ is $(\sum_{v\in V(G)}t_v)/3$.

Clearly, there are $2t_v$ vertices in $N_G(v)$ adjacent to $m-1$ vertices in $N_{\overline{G}}(v)$ and $m+1-2t_v$ vertices in $N_G(v)$ adjacent to $m$ vertices in $N_{\overline{G}}(v)$.

Since $0\le t_v\le\lfloor (m+1)/2\rfloor=m/2$, and each vertex in $N_G(v)$ is adjacent to at most one other vertex in $N_G(v)$,
the number of edges with one endpoint in $N_G(v)$ and the other in $N_{\overline{G}}(v)$ is $2t_v(m-1)+(m+1-2t_v)m$. Since $|N_{\overline{G}}(v)|=m^2+2$, it follows that $2t_v(m-1)+(m+1-2t_v)m\le m^2+2$ and hence $t_v\ge m/2-1$.

Assuming that $t_v=m/2$ for every $v\in V(G)$,  the number of triangles in $G$ would be $\left(\sum_{v\in V(G)}t_v\right)\!/3=m(m^2+m+4)/6.$ Since $m\equiv 2 \,(\text{mod}\, 6)$, let $a=(m-2)/6$, which is an integer. The total number of triangles then becomes $36a^3+42a^2+20a+10/3$, which is not an integer, leading to a contradiction. Therefore, there must exist a vertex $v_0\in V(G)$ such that $t_{v_0}=m/2-1$.

Let $F=G[V(G)\setminus(N_G(v_0)\cup\{v_0\})]$, so $|V(F)|=m^2+2$. Let $u_1,\ldots,u_{m+1}$ be the vertices in $N_G(v_0)$, with $u_{2i-1}$ adjacent to $u_{2i}$ for $1\le i\le m/2-1$. Define $A_j=N_G(u_j)\cap V(F)$ for $1\le j\le m+1$. Then, $|A_j|=m-1$ for $1\le j\le m-2$ and $|A_j|=m$ for $m-1\le j\le m+1$. Since $G\nsupseteq C_4$, if $i\ne j$, then $A_i\cap A_j=\emptyset$. Thus $|\bigcup_{i=1}^{m+1}A_i|=\sum_{i=1}^{m+1}|A_i|=(m-2)(m-1)+3m=m^2+2$, implying $\bigcup_{i=1}^{m+1}A_i=V(F)$.

Since $|A_1|=m-1$ is odd and $\delta(F[A_1])\le 1$, there exists $w_1\in A_1$ that is not adjacent to any other vertex in $A_1$. If $w_1$ is adjacent to two vertices in $A_i$, for some $i$, those two vertices  together with $w_1$ and $u_i$ would form a $C_4$, leading to $w_1$ being adjacent to at most one vertex in $A_i$.

If $w_1$ were adjacent to a vertex in $A_2$, then, since $u_1$ and $u_2$ are adjacent, a $C_4$ would be formed. Thus, $w_1$ is adjacent to $u_1$ and at most one vertex in each $A_i$ with $3\le i\le m+1$. Consequently, $w_1$ is adjacent to at most $m$ vertices in $G$, which leads to a contradiction and yields the desired result.
\end{proof}

\end{document}
